\section{Benchmarking Real-World Agents}

\subsection{Benchmark Pipeline}

Anthropic descompone la capacidad del Agent en varias métricas: autonomía, invocación de herramientas, profundidad de razonamiento, consistencia de preferencias. Cada ronda de evaluación encadena las siguientes etapas:

\begin{enumerate}
    \item \textbf{Design scenarios}: selección de tareas de ciclo largo (investigación, planeación, ejecución, reporte)
    \item \textbf{Establish baselines}: comparación con equipos humanos o modelos de generaciones anteriores
    \item \textbf{Instrument data}: recopilación de los registros de acciones del agent, las llamadas a la API y el trail de auditoría
    \item \textbf{Score metrics}: calificación automatizada + revisión humana
    \item \textbf{Feedback to training}: cada failure mode detectado se convierte, uno por uno, en un dataset/behavioral trace
\end{enumerate}

\begin{figure}[H]
\centering
\includegraphics[page=3,width=0.92\textwidth]{slides.pdf}
\caption{Agent Benchmark pipeline illustrated on slide 3.\protect\footnotemark}
\end{figure}
\footnotetext{Slides emphasize multi-step evaluation loops.}

\subsection{Case Study: Decision Support Task}

Una tarea típica consiste en que el Agent lea un documento interno de estrategia de 20 páginas, genere un resumen POAP (Plan, Options, Actions, Priority) y programe una reunión de follow-up.

\begin{itemize}
    \item El Agent necesita mantener la coherence entre páginas (anidamiento de párrafos, referencias)
    \item La salida debe atender a la vez los hechos, la cadena de razonamiento y las recomendaciones de acción
    \item Al final, un reviewer humano asigna la calificación según la template de `POAP`
\end{itemize}

\begin{knowledgebox}{Tailored judgment rubric}
Anthropic elabora un rubric para cada tipo de tarea: exactitud factual, integridad estructural, cumplimiento en la invocación de herramientas, si ofrece o no recomendaciones accionables. Solo se avanza a la actualización de ASL cuando todas las dimensiones son satisfactorias.
\end{knowledgebox}

\subsection{Resumen de la sección}

El benchmark pipeline de Anthropic encadena varios escenarios de tarea, la recopilación de datos, la calificación y la retroalimentación al entrenamiento en un ciclo cerrado; cada scenario cuenta con un rubric específico, con lo cual la evaluación de capacidades se convierte en una señal de gobernanza accionable.

\subsection{Evidence Matrix}

Para convertir la evaluación de capacidades en información para la toma de decisiones, Anthropic usa una \textbf{evidence matrix}: enumera, por etapa de la tarea, las métricas clave, las fuentes de datos, el desempeño actual, y da un límite de seguridad (green/orange/red).

\begin{table}[H]
\centering
\begin{tabular}{p{0.18\textwidth}p{0.24\textwidth}p{0.24\textwidth}p{0.24\textwidth}}
\toprule
Etapa & Métrica & Fuente de datos & Señal de riesgo actual \\
\midrule
Planning & clarity of plan, dependency awareness & prompt logs, human evaluation & plan misses key stakeholders \\
Tool invocation & tool success rate, error rate & API logs, retry counts & spike in `api\_error` events \\
Execution & prospect completion quality & manager reviews, outcome diff & 5\% of tasks reworked \\
Reporting & factual accuracy & summary diff + external documents & hallucinated reference detected \\
\bottomrule
\end{tabular}
\caption{Evidence matrix used for capability snapshots}
\end{table}

\begin{warningbox}{Evidence matrix refresh cadence}
Matrix rows are refreshed every week or after any significant dataset/model update. Without timely refresh, the safety gates rely on stale performance snapshots.
\end{warningbox}
